% ---
% title: Loss given default
% tags: finance, distribution-function, density
% ---

\begin{problem}
	\begin{context}
		When a borrower defaults, that is, stops paying back a loan, the bank
		usually recovers part of the money, for example by selling the
		collateral. The \emph{loss given default} (LGD) is the fraction of the
		loan that is lost. An LGD of 0 means that everything is recovered, an
		LGD of 1 means that nothing is.
	\end{context}
	The loss given default $L$ of a loan is a random number between 0 and 1
	with distribution function
	\[ F(x)=\begin{cases}
			0,   & x<0,          \\
			x^2, & 0\le x\le 1,  \\
			1,   & x>1.
		\end{cases} \]
	(a) Find the density $f(x)$ of $L$.
	(b) Find $P(L>0.5)$.
	(c) Find $P(0.2<L\le0.6)$.
	(d) Find the median loss: the number $m$ with $P(L\le m)=\frac12$.
\end{problem}

\begin{slideproblem}
	The loss given default $L$ of a loan has distribution function
	$F(x)=x^2$ for $0\le x\le1$ ($F(x)=0$ for $x<0$, $F(x)=1$ for $x>1$).
	(a) Find the density $f(x)$ of $L$.
	(b) Find $P(L>0.5)$.
	(c) Find $P(0.2<L\le0.6)$.
	(d) Find the median loss: the number $m$ with $P(L\le m)=\frac12$.
\end{slideproblem}

\begin{slides}
	\small
	\exstatement
	\pause
	\begin{columns}[T]
		\begin{column}{.62\textwidth}
			\sol (a) $f=F'$: \ $\alert{f(x)=2x}$ for $0\le x\le1$, \ $f(x)=0$
			otherwise.\\
			Check: $\int_0^1 2x\,dx=\big[x^2\big]_0^1=1$. Large losses are more
			likely.\\[1mm]
			\pause
			(b) $P(L>0.5)=1-F(0.5)=1-0.5^2=\alert{0.75}$.\\[1mm]
			\pause
			(c) $P(0.2<L\le0.6)=F(0.6)-F(0.2)=0.36-0.04=\alert{0.32}$.\\[1mm]
			\pause
			(d) $F(m)=m^2=\frac12$, \ so \ $m=\sqrt{0.5}\approx\alert{0.707}$:
			half of the defaults lose more than 70.7\% of the loan.
		\end{column}
		\begin{column}{.36\textwidth}
			\centering
			\begin{tikzpicture}[xscale=2.3,yscale=1.05]
				\draw[gray!40,very thin] (0,1)--(1.3,1) (0,2)--(1.3,2);
				\node[left,font=\tiny] at (0,1) {$1$};
				\node[left,font=\tiny] at (0,2) {$2$};
				\fill[teal!22] (.5,0) -- (.5,1) -- (1,2) -- (1,0) -- cycle;
				\draw[teal] (.5,0) -- (.5,1);
				\node[font=\tiny,text=teal!70!black] at (.78,.5) {$0.75$};
				\draw[->] (0,0) -- (1.4,0) node[right,font=\scriptsize] {$x$};
				\draw[->] (0,0) -- (0,2.35);
				\draw[navy,thick] (0,0) -- (1,2);
				\draw[navy,dashed] (1,0) -- (1,2);
				\draw[navy,thick] (1,0) -- (1.3,0);
				\draw[gold,thick] plot[domain=0:1,samples=50] (\x,{\x*\x}) -- (1.3,1);
				\node[font=\tiny,text=navy,left] at (.95,2.1) {$f(x)=2x$};
				\node[font=\tiny,text=gold!80!black,above] at (1.2,1) {$F(x)$};
				\foreach \x in {0.5,1}
				\node[below,font=\tiny] at (\x,0) {$\x$};
			\end{tikzpicture}\\
			{\tiny shaded area $=P(L>0.5)$}
		\end{column}
	\end{columns}
\end{slides}

\begin{solution}
	\idea The density of a continuous random variable is the derivative of its
	distribution function, $f=F'$. Probabilities of intervals come directly
	from $F$: $P(a<L\le b)=F(b)-F(a)$, and $P(L>a)=1-F(a)$.

	\step{(a)} Differentiate $F$ piece by piece. For $x<0$ and for $x>1$ the
	function $F$ is constant (0 and 1), so its derivative is 0 there. For
	$0<x<1$ we have $F(x)=x^2$, so $F'(x)=2x$. Hence
	\[ f(x)=\begin{cases}
			2x, & 0\le x\le1, \\
			0,  & \text{otherwise}.
		\end{cases} \]
	(The value of a density at the single points $x=0$ and $x=1$ does not
	matter: a single point has zero area under it.) Check that the total area
	under the density is 1. An antiderivative of $2x$ is $x^2$, so
	\[ \int_0^1 2x\,dx=\Big[x^2\Big]_0^1=1^2-0^2=1. \]
	The density grows with $x$: losses close to 1 are more likely than losses
	close to 0. So in this model, when this borrower defaults, the bank
	typically loses a large part of the loan.
	\begin{center}
		\begin{tikzpicture}[scale=2]
			\fill[teal!30] (.5,0) -- (.5,1) -- (1,2) -- (1,0) -- cycle;
			\draw[->] (-.1,0) -- (1.4,0) node[right] {\small $x$};
			\draw[->] (0,-.1) -- (0,2.25) node[above] {\small $y$};
			\draw[gold,thick] plot[domain=0:1,samples=50] (\x,{\x*\x});
			\draw[gold,thick] (1,1) -- (1.35,1);
			\draw[navy,thick] (0,0) -- (1,2);
			\draw[navy,dashed] (1,0) -- (1,2);
			\draw[navy,thick] (1,0) -- (1.35,0);
			\draw (.5,.03) -- (.5,-.03) node[below] {\scriptsize $0.5$};
			\draw (1,.03) -- (1,-.03) node[below] {\scriptsize $1$};
			\draw (.03,1) -- (-.03,1) node[left] {\scriptsize $1$};
			\draw (.03,2) -- (-.03,2) node[left] {\scriptsize $2$};
			\node[navy,left] at (.95,2.08) {\small $f(x)=2x$};
			\node[gold!80!black,right] at (1.05,1.12) {\small $F(x)$};
			\node[teal!60!black] at (.82,.28) {\small $0.75$};
		\end{tikzpicture}
	\end{center}
	The navy line is the density, the gold curve is the distribution
	function. The shaded area is $P(L>0.5)$, computed in (b).

	\step{(b)} The event $\{L>0.5\}$ is the complement of $\{L\le0.5\}$, and
	$P(L\le0.5)=F(0.5)$. So
	\[ P(L>0.5)=1-F(0.5)=1-0.5^2=1-0.25=0.75. \]
	The same result from the density: the shaded region in the picture is a
	trapezoid with parallel sides $f(0.5)=1$ and $f(1)=2$ and width $0.5$, so
	its area is $\frac{1+2}{2}\cdot0.5=0.75$. With an integral:
	$\int_{0.5}^1 2x\,dx=\big[x^2\big]_{0.5}^1=1-0.25=0.75$.

	\step{(c)} The event $\{L\le0.6\}$ splits into two disjoint parts,
	$\{L\le0.2\}$ and $\{0.2<L\le0.6\}$. So $F(0.6)=F(0.2)+P(0.2<L\le0.6)$,
	which gives the rule $P(a<L\le b)=F(b)-F(a)$:
	\[ P(0.2<L\le0.6)=F(0.6)-F(0.2)=0.6^2-0.2^2=0.36-0.04=0.32. \]
	For a continuous random variable it does not matter whether the endpoints
	are included: $P(L=0.2)=0$, so $P(0.2\le L\le0.6)$ is also $0.32$.

	\step{(d)} We need the point $m$ where $F$ reaches $\frac12$. Since $m$
	must lie between 0 and 1, here $F(m)=m^2$, and
	\[ m^2=\frac12 \quad\Longrightarrow\quad m=\sqrt{\frac12}=\sqrt{0.5}\approx0.707 \]
	(we take the positive root, because $m$ is between 0 and 1). So in half of
	the defaults the bank loses at most $70.7\%$ of the loan, and in the other
	half it loses more.
	\begin{answer}
		(a) $f(x)=2x$ for $0\le x\le1$, $f(x)=0$ otherwise
		\quad (b) $P(L>0.5)=0.75$
		\quad (c) $P(0.2<L\le0.6)=0.32$
		\quad (d) $m=\sqrt{0.5}\approx0.707$
	\end{answer}
	\begin{remark}
		\textbf{Common mistake:} answering (b) with $1-0.5=0.5$, as if $L$ were
		uniform on $[0,1]$. The distribution function is $x^2$, not $x$, so
		the probabilities must be read off $F$. Another mistake is to compute
		(c) as $F(0.6-0.2)=F(0.4)=0.16$: the probability of an interval is the
		\emph{difference} of $F$ at the two endpoints, not $F$ of the length.
	\end{remark}
\end{solution}
